\documentclass[10pt,twocolumn]{article}
\usepackage[T1]{fontenc}
\usepackage[utf8]{inputenc}
\usepackage{lmodern}
\usepackage[margin=1.8cm,top=2.4cm,bottom=2cm,columnsep=0.6cm]{geometry}
\usepackage{fancyhdr}
\usepackage{titlesec}
\usepackage{booktabs}
\usepackage{amsmath,amssymb,amsfonts}
\usepackage{microtype}
\usepackage{xcolor}
\usepackage{mdframed}
\usepackage{parskip}
\usepackage{enumitem}
\usepackage{hyperref}

\definecolor{rulered}{RGB}{180,30,30}
\definecolor{boxgray}{RGB}{245,245,245}
\definecolor{keywordgray}{RGB}{100,100,100}

\pagestyle{fancy}
\fancyhf{}
\fancyhead[L]{\textsc{\small Claude Mag}}
\fancyhead[C]{\small\textit{Machine Learning Research Digest}}
\fancyhead[R]{\small 2026-10-19}
\fancyfoot[C]{\small\thepage}
\renewcommand{\headrulewidth}{0.8pt}

\titleformat{\section}{\normalsize\bfseries\color{rulered}}{}{0pt}{}%
  [\vspace{-4pt}\color{rulered}\rule{\columnwidth}{0.4pt}]
\titlespacing{\section}{0pt}{8pt}{4pt}

\hypersetup{colorlinks=true,urlcolor=rulered,linkcolor=black}

\begin{document}
\setlength{\parindent}{0pt}
\setlength{\parskip}{4pt}

{\small\color{keywordgray}\textbf{Keywords:}
  video benchmark \textbullet{}
  omni-modal \textbullet{}
  temporal grounding \textbullet{}
  audio-visual understanding}\par\vspace{4pt}

{\LARGE\bfseries\raggedright
  OmniCompass: A Holistic Benchmark for Omni-Modal Video Grounding
  and Understanding}\par\vspace{4pt}

{\large\itshape\raggedright
  Paired-Setting Evaluation of 50 MLLMs Reveals Widespread
  Pattern-Matching and Persistent Grounding Failures in Both
  Vision-Language and Audio-Visual Tracks}\par\vspace{6pt}

{\small\textbf{By} Xuyang Liu, Shihao Wang, Xinyu Wei, Minghan Li,
  Zeqing Wang, Zuhao Yang, Zichen Wen \& Lei Zhang}\par\vspace{2pt}

{\small\color{keywordgray}arXiv:2610.18829 \textbullet{} October 2026}
\par\vspace{2pt}\noindent{\color{rulered}\rule{\columnwidth}{0.6pt}}\vspace{6pt}

Today's video benchmarks test language or vision in isolation, leaving
the joint audio-visual grounding capabilities of omni-modal models
unmeasured. \emph{OmniCompass} introduces a paired evaluation
protocol---each instance evaluated under text-only \emph{and}
text$+$evidence conditions---across two tracks: OmniCompass-VL
(820 videos, 2,460 instances) and OmniCompass-AV (240 videos, 720
instances from narrative-rich, audio-dense content). Evaluation of 50
MLLMs finds that open-source models trail proprietary ones by 15--20
points on VL tasks and more on AV tasks; widespread paired-setting
inconsistency ($>$30\% of instances) exposes surface-level pattern
matching in even the best-performing models.

\begin{mdframed}[backgroundcolor=boxgray,linecolor=rulered,linewidth=1pt,
    innerleftmargin=6pt,innerrightmargin=6pt,
    innertopmargin=6pt,innerbottommargin=6pt]
\textbf{\small AT A GLANCE}\par\vspace{4pt}
\begin{description}[leftmargin=0pt,labelsep=4pt,itemsep=2pt,parsep=0pt,
    font=\small\bfseries]
  \item[Problem:] Video benchmarks test language or vision in isolation;
    no systematic measurement of joint audio-visual grounding exists.
  \item[Why it matters:] Omni-modal models are deployed in video
    understanding; benchmarks must reveal whether they genuinely ground
    queries in evidence or pattern-match instruction formats.
  \item[Method:] Two-track benchmark with paired evaluation
    (text-only/text$+$evidence); 8 VL task categories and 5 AV
    categories; temporal grounding, spatial grounding, and QA subtasks.
  \item[Result:] 50 MLLMs tested; open-source trails proprietary by
    15--20 pp on VL and more on AV; $>$30\% paired-setting inconsistency
    in the majority of models.
  \item[Evidence:] 820$+$240 videos, 2{,}460$+$720 instances;
    inter-annotator $\kappa > 0.8$.
\end{description}
\end{mdframed}

\section{The Big Picture}

Omni-modal models can now process video, audio, and text simultaneously,
yet their evaluation lags: benchmarks still test modalities in isolation.
A model that answers a temporal grounding question correctly under
text-only conditions may fail when the same question is accompanied by
contradictory visual evidence---or vice versa. This mismatch is invisible
to single-condition evaluation.

OmniCompass is built around a paired evaluation protocol that tests every
instance twice: once with text alone and once with the relevant visual or
audio evidence. A model that genuinely understands the target should
respond consistently (and correctly) under both conditions. Inconsistency
across conditions is a diagnostic signal: the model is pattern-matching
the instruction format rather than grounding in the evidence.

\section{Why Existing Approaches Fall Short}

\textbf{Standard VQA benchmarks} test single-image comprehension; they
do not cover video-level temporal or spatial grounding across a narrative.

\textbf{Text-only LLM benchmarks} ignore visual and audio modalities
entirely; they cannot measure multimodal grounding by construction.

\textbf{Existing video QA benchmarks} (MVBench, VideoMME, etc.) focus on
visual understanding without audio-dependent tasks or the paired-setting
probe that would reveal instruction-format reliance.

\textbf{Prior audio-visual benchmarks} are small-scale and lack the
structured temporal/spatial annotation and the paired consistency
diagnostic that OmniCompass provides.

\section{Core Mechanism: Paired Evaluation Design}

\textbf{Two evaluation tracks:}

\emph{OmniCompass-VL} (Vision-Language): 820 videos, 2{,}460 instances.
Questions require visual evidence. Eight task categories: five perception
(object recognition, action recognition, attribute identification,
counting, OCR) and three reasoning (causal, commonsense, temporal
reasoning). Three subtask types: temporal grounding (TG), spatial
grounding (SG), and question answering (QA).

\emph{OmniCompass-AV} (Audio-Visual): 240 videos, 720 instances from
narrative-rich, audio-dense content. Questions require joint audio-visual
evidence. Five task categories: three perception (speech recognition,
sound identification, speaker localization) and two reasoning
(audio-visual causal reasoning, event coordination). Same TG/SG/QA
subtask structure.

\textbf{Paired evaluation:} Each instance is evaluated twice per model:
text-only and text$+$evidence. Inconsistency between conditions flags
surface-level pattern matching.

\section{Technical Formulation}

For model $\mathcal{M}$, query $q_i$, visual/audio evidence $v_i$, and
ground-truth answer $a_i^*$, define per-instance scores:
\[
s_i^{\text{text}} = \mathbf{1}\!\left[\mathcal{M}(q_i) = a_i^*\right],\quad
s_i^{\text{multi}} = \mathbf{1}\!\left[\mathcal{M}(q_i, v_i) = a_i^*\right]
\]

The \textbf{consistency score} across $N$ instances is:
\[
C(\mathcal{M}) = \frac{1}{N}\sum_{i=1}^{N}
  \mathbf{1}\!\left[s_i^{\text{text}} = s_i^{\text{multi}}\right]
\]

A high $C(\mathcal{M})$ with high $\sum_i s_i^{\text{multi}}$ indicates
reliable grounding; a high $C(\mathcal{M})$ with low accuracy indicates
consistent failure; a low $C(\mathcal{M})$ regardless of accuracy
indicates format sensitivity.

Grounding precision is measured separately for temporal and spatial tasks.
For a predicted interval $[\hat{t}_s, \hat{t}_e]$ and ground-truth
$[t_s, t_e]$:
\[
\mathrm{T\text{-}IoU} =
\frac{[\hat{t}_s, \hat{t}_e]\cap[t_s, t_e]}{[\hat{t}_s, \hat{t}_e]\cup[t_s, t_e]}
\]
Spatial grounding uses standard bounding-box IoU.

\section{Learning or Inference Procedure}

OmniCompass is a benchmark; no model training is involved. Evaluation
protocol:

\begin{enumerate}[itemsep=2pt,parsep=0pt]
  \item Collect videos with expert-annotated temporal segments, spatial
    bounding boxes, and QA pairs (inter-annotator $\kappa > 0.8$).
  \item Adversarially filter questions solvable from text alone.
  \item Construct paired instances: text-only and text$+$evidence versions
    of every question.
  \item Evaluate each model in zero-shot mode on all instances.
  \item Compute per-track accuracy, T-IoU, spatial IoU, and consistency
    scores; stratify by task category and model family.
\end{enumerate}

\section{What the Guarantee Says}

\begin{itemize}[itemsep=2pt,parsep=0pt]
  \item Majority of evaluated models show \textbf{consistency scores
    below 70\%}: over 30\% of instances produce different correctness
    outcomes depending on whether evidence is provided.
  \item Grounding precision (temporal IoU, spatial IoU) remains
    \textbf{low across all model families}, including proprietary ones.
  \item Open-source models \textbf{trail proprietary by 15--20 pp} on
    the VL track; the gap widens on the AV track.
  \item Models that score well on VL tasks \textbf{do not uniformly
    transfer} to AV tasks, showing audio-visual reasoning requires more
    than visual competence.
\end{itemize}

\section{Experimental Findings}

\begin{itemize}[itemsep=2pt,parsep=0pt]
  \item \textbf{Models (50 total):} Proprietary: GPT-4o, Claude Opus,
    Gemini models. Open-source: Qwen-VL, LLaVA variants, InternVL.
  \item \textbf{VL accuracy gap:} Proprietary average 15--20 pp above
    open-source on QA subtasks; smaller gap on perception, larger on
    reasoning.
  \item \textbf{AV accuracy gap:} The proprietary-to-open-source gap
    widens; several open-source models perform near chance on
    audio-visual causal reasoning.
  \item \textbf{Consistency:} Majority of models below 70\%; temporal
    reasoning and causal reasoning show the lowest consistency.
  \item \textbf{Grounding:} T-IoU and spatial IoU are low across all
    families, confirming grounding---not just QA---is the bottleneck.
\end{itemize}

\section{Ablations and Interpretation}

\textbf{Paired-gap analysis:} Models with the smallest text-vs-multimodal
accuracy gap tend to be the most accurate overall; genuine grounding
correlates with benchmark performance.

\textbf{Category breakdown:} Temporal reasoning and causal reasoning show
the lowest consistency scores, suggesting models rely heavily on
surface-level linguistic cues for these harder categories.

\textbf{Audio-visual transfer:} Models that excel at VL tasks do not
uniformly extend their advantage to AV tasks; audio-visual reasoning
requires explicitly joint processing that visual-only training does not
provide.

\textbf{Adversarial filtering:} Questions solvable from text alone
inflate apparent model capability; after filtering, accuracy on the
remaining paired set drops substantially, exposing true grounding rates.

\textbf{Annotation quality:} All task categories meet inter-annotator
agreement $\kappa > 0.8$; the consistency-score diagnostic is robust to
annotation noise because it compares a model against itself across
conditions, not against potentially noisy labels.

\section*{Reference}

Xuyang Liu, Shihao Wang, Xinyu Wei, Minghan Li, Zeqing Wang, Zuhao Yang,
Zichen Wen, Lei Zhang.
\textbf{OmniCompass: A Holistic Benchmark for Omni-Modal Video Grounding
and Understanding}.
arXiv:2610.18829, October 2026.\\
\url{https://arxiv.org/abs/2610.18829}

\end{document}
